\chapter{Literature Review}

Sleep plays a vital role in maintaining overall health, and its assessment depends on the accurate classification of sleep stages. This classification is crucial for diagnosing sleep disorders, including insomnia, sleep apnea, and narcolepsy. Traditionally, sleep experts analyze Polysomnography (PSG) data, which captures multiple physiological signals such as Electroencephalography (EEG), Electrocardiography (ECG), Electrooculography (EOG), and Electromyography (EMG). However, manual sleep stage classification is time-consuming, prone to inter-rater variability, and requires specialized expertise. As a result, researchers have increasingly focused on deep learning-based approaches to automate and enhance sleep stage classification. \\

Deep learning models, particularly Convolutional neural networks (CNNs) and Recurrent neural networks (RNNs), have demonstrated strong performance in sleep analysis by capturing complex patterns in PSG data. However, CNNs primarily extract local features, limiting their ability to model long-range dependencies, while RNNs struggle with capturing long-term temporal dependencies in sleep signals. Recently, vision transformers (ViTs) have gained attention due to their ability to model global contextual relationships using self-attention mechanisms. Unlike traditional deep learning models, ViTs can efficiently process PSG signals transformed into image representations, introducing a novel approach to sleep stage classification. \\

This literature review explores key research on deep learning-based sleep stage classification, with a particular focus on the role of ViTs. It examines various methods for converting PSG signals into images and compares traditional data augmentation techniques with Generative adversarial networks (GANs) for improving model accuracy. Additionally, the review highlights critical challenges, including class imbalance and inter-subject variability. By analyzing existing studies, this work aims to provide insights into the advantages and limitations of using ViTs for sleep stage classification.
%--------------------------------------------------------------------------------------------------------

\section{Traditional Approaches to Sleep Stage Classification}
\label{lit:1}

Traditional approaches have largely relied on machine learning based methods, involving handcrafted feature extraction followed by classical classification algorithms. These methods have contributed significantly to sleep staging research. However, they exhibit limitations in generalizability and computational efficiency, leading to a transition toward deep learning-based techniques.

\subsection{Machine Learning Based Approaches}

Feature engineering has been a fundamental step in traditional sleep stage classification. Studies have demonstrated the effectiveness of extracting time-domain, frequency-domain, and nonlinear features from Polysomnography (PSG) signals. Frequency-domain analysis, particularly \ac{PSD}, has been instrumental in distinguishing sleep stages based on spectral band activity \parencite{prerau2017sleep}. Additionally, nonlinear features, such as approximate entropy, have been employed to quantify the complexity of \acs{PSG} signals \parencite{acharya2018deep}. However, despite their effectiveness, handcrafted features require expert knowledge and often lack robustness across diverse datasets. \\

Several classical machine learning algorithms have been applied to sleep stage classification. \acp{SVM}  have been widely adopted due to their strong theoretical foundation and competitive performance \parencite{lajnef2015learning}. However, they suffer from scalability issues and sensitivity to class imbalance. Random Forests have shown promise in handling nonlinearity in \acs{PSG} data, though they require extensive hyperparameter tuning \parencite{liu2020automatic}. Similarly, \ac{KNN} has been explored for its simplicity but struggles with high-dimensional feature spaces \parencite{qureshi2018human}. \ac{HMM} have been employed to model the sequential nature of sleep stages, leveraging transition probabilities \parencite{doroshenkov2007classification}. However, \acp{HMM} assume Gaussian distributions, which may not accurately capture the complexity of \acs{PSG} signals. While these models have demonstrated moderate success, their reliance on feature engineering and limited adaptability across subjects have driven researchers toward deep learning-based approaches.

\subsection{Early Deep Learning Approaches}

With advancements in deep learning, Convolutional Neural Network (CNN) have emerged as a powerful tool for automatic feature extraction from raw Polysomnography (PSG) data. Studies have shown that 1D \acp{CNN} effectively capture spatial dependencies in Electroencephalography (EEG) signals, eliminating the need for manual feature engineering. Furthermore, 2D \acp{CNN}, which transform \acs{EEG} signals into spectrograms, have demonstrated superior performance by leveraging time-frequency representations \parencite{biswal2018expert}. However, \acp{CNN} have limitations, particularly in modeling long-range temporal dependencies \parencite{phan2018joint}, which restricts their ability to fully capture sleep stage transitions. \\

To address the temporal modeling limitations of \acp{CNN}, \acp{RNN} and their variant, \ac{LSTM} networks, have been adopted for sleep stage classification.  Studies have shown that \acp{LSTM} mitigate vanishing gradient issues and effectively capture long-term dependencies in \acs{PSG} sequences \parencite{supratak2017deepsleepnet}. Furthermore, hybrid models combining \acp{CNN} for feature extraction and \acp{LSTM} for sequence modeling have achieved improved performance \parencite{chambon2018deep}. Despite these advances, \acp{RNN} and \acp{LSTM} require significant computational resources and may struggle with long \acs{PSG} recordings, leading researchers to explore attention-based mechanisms and Transformer models \parencite{vaswani2017attention}.

%--------------------------------------------------------------------------------------------------------
\section{Transformers in Sleep Stage Classification}
\label{lit:2}

Transformers, originally designed for \ac{NLP}, have proven effective in medical signal analysis, including sleep stage classification. Their ability to capture long-range dependencies makes them well suited for analyzing complex Polysomnography (PSG) signals. This section reviews transformer-based approaches for sleep stage classification, compares them with other deep learning models.

\subsection{Transformers in Medical Signal Analysis}

Several studies have demonstrated the effectiveness of transformers in medical signal processing, particularly for Electroencephalography (EEG) and Electrocardiography (ECG) classification. Studies have used  transformer-based models to EEG classification and found that self-attention mechanisms improved the results compared to Convolutional Neural Networks (CNNs) and Recurrent Neural Networks (RNNs) \parencite{UnknownAuthor2024}. Similarly, studies have used transformer models for ECG classification, demonstrating that transformers outperform CNNs in recognizing abnormal heart rhythms due to their superior ability to capture long-term dependencies in time series data \parencite{krishna2024enhanced}. These findings suggest that transformers, with their ability to extract both local and global features, are well suited for analyzing physiological signals, including PSG recordings used in sleep staging. \\

In the context of sleep stage classification, studies have shown that transformers can achieve higher accuracy and robustness than CNNs. Researchers have conducted a comparative study on EEG based sleep staging, where a transformer model outperformed a CNN in distinguishing Rapid Eye Movement (REM) sleep and deep sleep stages \parencite{thuwajit2021eegwavenet}. The authors attributed this improvement to the self-attention mechanism, which allowed the transformer to consider long-term dependencies across sleep cycles rather than focusing only on local spatial features, as CNNs do. Similarly, studies have proposed a hybrid CNN-transformer model that leveraged the feature extraction power of CNNs alongside the contextual learning ability of transformers \parencite{yao2023cnn}. Their results indicated that adding transformer layers improved classification performance, particularly in recognizing transitional sleep stages, which are often misclassified by CNN-based models. \\

Further supporting these findings, A group of researchers have trained a Vision Transformer (ViT) model on spectrogram representations of PSG signals and reported higher F1-scores and Cohen’s kappa values than conventional CNN models \parencite{kurniawanvision}. The ViT-based approach was particularly effective in capturing subtle temporal variations in sleep signals, which CNNs often struggle with due to their reliance on fixed kernel sizes. Moreover, studies have highlighted that transformers perform better than CNNs in cases with high inter-subject variability, a common challenge in sleep stage classification \parencite{guo2024flexsleeptransformer}. While CNNs excel at extracting local spatial patterns, they often fail to generalize well across different individuals due to variations in sleep physiology. In contrast, transformers demonstrate greater adaptability by capturing both global and local dependencies, making them more robust across diverse sleep datasets.\\

ViTs offer several key advantages that make them particularly suited for sleep stage classification. First, they excel at capturing long-range dependencies, which is crucial for sleep analysis since sleep stages transition over extended periods. Unlike CNNs, which rely on local receptive fields, ViTs process entire input sequences simultaneously, allowing them to recognize patterns across full sleep cycles \parencite{kazemi2025multimodal}. This capability helps in distinguishing between sleep stages that share similar spectral characteristics but occur at different times in the sleep process. \\

Self-attention mechanisms enable ViTs to assign different levels of importance to various regions of an input image, making them more effective at identifying critical sleep-related features \parencite{lee2025explainable}. For example, in spectrogram-based PSG analysis, ViTs can focus on frequency components most relevant to specific sleep stages, rather than treating all input regions equally. This contrasts with CNNs, which apply convolutional filters uniformly across an image, potentially overlooking subtle but important features. Another major advantage of ViTs is their flexibility in handling multi-modal data. Sleep staging often requires integrating multiple signals, such as EEG, ECG, and EOG, to achieve high classification accuracy. Traditional CNNs are typically designed for single-modal inputs, requiring manual feature fusion techniques to combine information from different signals. In contrast, ViTs can process multi-modal PSG data in a unified framework, extracting meaningful representations from diverse physiological signals without explicit feature engineering \parencite{kazemi2025multimodal}.\\

These studies collectively suggest that transformers, particularly ViT, offer a significant advantage over CNN-based models for sleep stage classification. Their ability to process entire sleep cycles holistically, rather than relying solely on local patterns, allows for more accurate and reliable sleep staging.  \\

%--------------------------------------------------------------------------------------------------------

\section{Transforming PSG Signals into Image Representations}
\label{lit:3}

PSG signals, including EEG, EOG, EMG, and ECG, are traditionally analyzed as time series data. However, with the emergence of ViTs, which require image-based inputs, an essential step in sleep stage classification is the transformation of these sequential signals into meaningful image representations. This section explores various transformation techniques, including time-frequency methods, recurrence plots, and other advanced encoding strategies, analyzing their impact on model performance.

\subsection{Time-Frequency Transformations}

Time-frequency transformations are essential for analyzing physiological signals like EEG and PSG, enabling effective sleep stage classification. Techniques such as \ac{STFT}, \ac{CWT}, and \ac{HHT} each offer unique advantages and limitations in capturing time-dependent frequency variations. Their effectiveness depends on signal characteristics and application needs. \\

STFT, one of the earliest time-frequency analysis techniques, partitions signals into overlapping windows and applies the Fourier Transform to each segment, producing a spectrogram that visualizes time, frequency, and power intensity \parencite{cohen2020time}. This method is particularly well suited to detect periodic patterns in EEG and PSG signals. However, a fundamental limitation of STFT lies in its reliance on fixed window sizes, leading to a resolution trade-off enhancing time resolution diminishes frequency resolution and vice versa \parencite{ali2024detection}. While this trade-off restricts its ability to analyze transient patterns effectively, STFT remains widely used due to its computational efficiency and straightforward implementation. Nevertheless, its rigidity in resolution makes it less effective for analyzing non-stationary and complex physiological signals, raising questions about its suitability in dynamic sleep stage classification scenarios. \\

In contrast, CWT overcomes STFT resolution trade-off by utilizing wavelets localized functions capable of simultaneously capturing time and frequency features at multiple scales \parencite{torrence1998practical}. This capability allows CWT to preserve fine-grained temporal features, making it particularly effective in detecting transient sleep patterns such as sleep spindles, K-complexes, and slow-wave sleep. Despite these advantages, CWT introduces computational challenges due to its increased complexity. Studies have  demonstrated that CWT-based scalograms yield superior classification performance compared to traditional spectrograms \parencite{schirrmeister2017deep}. However, its computational intensity limits its feasibility in real-time applications. While its adaptability makes it highly suitable for complex physiological signal analysis, the trade-off between accuracy and efficiency remains an important consideration. \\

HHT presents a more flexible alternative by directly adapting to the non-stationary nature of physiological signals. Through empirical mode decomposition, HHT decomposes signals into intrinsic mode functions, followed by the Hilbert Transform to obtain instantaneous frequency representations \parencite{huang1998empirical}. This adaptability enables HHT to extract biologically meaningful features, making it particularly advantageous for analyzing EEG and PSG data. However, the method's high computational complexity poses a significant challenge, particularly in large-scale applications. While its capability to capture dynamic signal variations surpasses both STFT and CWT, its implementation requires substantial computational resources, limiting its practicality in real-time analysis.

\subsection{Gramian Angular Fields and Markov Transition Fields}

Time-series encoding techniques, such as \ac{GAF} and \ac{MTF}, have emerged as promising approaches for transforming PSG signals into structured representations suitable for deep learning-based classification. These methods offer complementary perspectives on temporal evolution, enhancing both interpretability and model performance. \\

GAF encodes time-series data into angular space, representing temporal evolution as a two-dimensional image \parencite{wang2015imaging}. This method preserves both magnitude and directional information, making it particularly useful for visualizing sleep stage transitions. However, GAF is highly sensitive to normalization techniques, as preprocessing variations can significantly alter encoded patterns, thereby affecting classification outcomes. This dependency raises concerns regarding the robustness of GAF-based transformations across diverse datasets and sleep environments. \\

MTF, in contrast, models probabilistic state transitions within a time series, capturing the evolving dynamics of sleep stages \parencite{wang2015imaging}. It has proven particularly effective in distinguishing REM from NREM sleep, as it directly encodes transition probabilities between states. Nevertheless, MTF introduces substantial computational overhead due to the need for Markov matrix calculations. This limitation reduces its suitability for real-time applications, despite its strong potential for capturing sleep architecture patterns. \\

Recent research suggests that hybridizing GAF and MTF enhances classification accuracy in deep learning architectures, particularly in ViTs \parencite{wang2015encoding}. The combination leverages GAF’s ability to encode fine-grained temporal structures and MTF’s capacity to capture state transitions, leading to a more comprehensive representation of sleep signals. While this hybrid approach appears promising, its practical feasibility remains an open question, given the increased computational complexity involved.

%--------------------------------------------------------------------------------------------------------

\section{Generative Adversarial Networks for Data Augmentation}
\label{lit:4}

Class imbalance in sleep stage classification affects model performance, as rare sleep stages (e.g., N3, REM) are often misclassified due to their lower frequency. Traditional augmentation techniques like time warping, noise addition, and window slicing help mitigate this but may not generate sufficiently diverse samples. GAN provide a more effective solution by creating synthetic PSG data that closely resembles real signals, improving model generalization. This section reviews the impact of class imbalance, examines studies using GANs for data augmentation, and compares traditional augmentation with GAN-based approaches.

\subsection{Class Imbalance in Sleep Stage Classification}

Sleep datasets often have an uneven distribution of sleep stages, with wake (W) and light sleep (N2) being more common than deep sleep (N3) and REM sleep \parencite{pardamean2023supervised}. This imbalance leads deep learning models to misclassify rare stages, as they tend to prioritize the majority classes during training. Models trained on imbalanced datasets typically exhibit lower recall and F1-scores for minority classes, which reduces their overall reliability in clinical applications. Researchers have attempted to mitigate this issue using oversampling, undersampling, and cost-sensitive learning, but these approaches may either introduce redundancy or remove valuable data \parencite{weiss2007cost}.

\subsection{GAN-Based vs Traditional Approaches}

Traditional data augmentation techniques, such as time shifting, noise addition, and flipping, have been widely used to increase the diversity of sleep datasets \parencite{ling2022staging}. While these methods are effective in expanding training data, they do not generate entirely new samples but instead modify existing ones. As a result, they may fail to introduce meaningful variations that could help deep learning models generalize better, particularly for underrepresented sleep stages. GAN offer a more advanced approach by learning the underlying data distribution and generating synthetic PSG signals that resemble real ones. \\

Several studies have demonstrated the benefits of using GANs for sleep stage classification. A team have trained a \ac{DCGAN} to generate synthetic EEG data and found that incorporating GAN-augmented samples significantly improved classification performance, particularly for rare sleep stages \parencite{du2024data}. Similarly, studies have used a \ac{cGAN} to generate class-specific EEG spectrograms, ensuring that synthetic samples preserved essential sleep-related features \parencite{fu2021conditional}. Their results showed that models trained with GAN-augmented data achieved higher accuracy, F1-score, and Cohen’s kappa than those trained on real data alone. \\

GANs also help address inter-subject variability, which is another challenge in sleep stage classification. Studies have proposed SleepGAN, a model designed to generate realistic EEG signals that capture diverse sleep patterns \parencite{cheng2024sleepegan}. Their study showed that training deep learning models with GAN-augmented data improved their ability to generalize across different subjects, reducing overfitting and increasing robustness. However, despite their advantages, GANs require careful training to prevent the generation of unrealistic or biased synthetic samples. Poorly trained GANs may produce artifacts or mode collapse, which can negatively impact classification performance. \\

Comparing traditional augmentation techniques with GAN-based approaches, studies indicate that GANs outperform conventional methods in enhancing sleep stage classification, particularly for rare sleep stages. Traditional augmentation methods are easier to implement and computationally less expensive but may not introduce sufficient variability in the dataset. In contrast, GANs generate entirely new synthetic data, capturing complex sleep patterns and improving model performance. Nonetheless, GAN-based augmentation requires extensive computational resources and careful tuning to ensure the generated samples retain physiological validity. Future research should focus on optimizing GAN architectures and integrating them with hybrid augmentation techniques to further improve the reliability of sleep stage classification models.

%--------------------------------------------------------------------------------------------------------
\section{Research Gaps and Unexplored Areas}
\label{lit:5}

Despite significant advancements in sleep stage classification, several critical gaps remain in the literature. While previous studies have explored traditional machine learning and deep learning techniques (Section \ref{lit:1}), as well as the application of Transformers to medical signal analysis (Section \ref{lit:2}), there has been little focus on optimizing pre-trained Vision Transformers (ViTs) for this task. Additionally, while various image transformation techniques have been investigated (Section \ref{lit:3}) and GAN-based augmentation strategies have been proposed (Section \ref{lit:4}), their impact on Transformer-based models remains largely unexplored. The following gaps highlight these limitations.

\subsection{Limited Use of Pre-Trained Vision Transformers in Sleep Stage Classification}
\label{lit:51}

As discussed in Section \ref{lit:2}, Transformers have shown significant promise in medical signal processing. However, most existing studies in sleep stage classification train these models from scratch, which often demands large annotated datasets and substantial computational resources. This approach poses a major limitation in the context of sleep research, where datasets are typically small, imbalanced, and difficult to obtain due to the complexity and cost of polysomnographic recordings. \\

Pre-trained Vision Transformers (ViTs), which have revolutionized tasks in natural image domains, remain underutilized in PSG-based sleep stage classification. Their ability to transfer learned visual representations to new domains could offer substantial advantages—such as improved generalization across subjects and conditions, faster convergence, and reduced risk of overfitting in low-data regimes. This is especially important in clinical applications, where patient variability and limited labeled data can hinder model reliability. \\

Despite these advantages, few studies have explored how pre-trained ViTs can be effectively fine-tuned for sleep stage classification. Understanding their performance in this context could lead to more efficient and accurate models, while also enabling the development of practical, real-time systems for automated sleep monitoring.

\subsection{Underexplored Comparisons Between Different Image Representation Techniques for PSG Data}
\label{lit:52}

Transforming time-series physiological signals—such as EEG, EOG, and EMG into visual formats is a crucial step when applying Vision Transformers (ViTs) to sleep stage classification. Several transformation techniques have been proposed in recent literature, including spectrograms, scalograms, Gramian Angular Fields (GAF), and Markov Transition Fields (MTF), each aiming to capture distinct temporal and frequency characteristics of sleep-related signals (Section \ref{lit:3}). \\

However, these studies often apply a single transformation method without systematically comparing the efficacy of different representations on Transformer-based models. This lack of comparative analysis creates uncertainty around which transformation best aligns with the spatial attention mechanisms of ViTs. Moreover, the suitability of each representation may vary depending on the sleep stage, modality (e.g., EEG vs. EMG), or signal resolution, further complicating model design. \\

A comprehensive evaluation of image representations is essential to identify which ones optimize feature richness, spatial consistency, and model interpretability for ViTs. Without such a comparison, model performance may be hindered by suboptimal input encoding, limiting the potential of even state-of-the-art Transformer architectures in sleep stage classification.

\subsection{Insufficient Use of GAN-Based Data Augmentation for Transformer Models}
\label{lit:53}

Class imbalance is a persistent challenge in sleep stage classification, particularly for underrepresented stages like N1 or REM, which often lead to skewed model performance (Section \ref{lit:4}). Generative Adversarial Networks (GANs) have emerged as powerful tools for addressing this imbalance by synthesizing realistic physiological data or image representations to augment training datasets. While previous studies have explored GAN-based augmentation for CNN-based sleep models, their integration into Transformer pipelines remains largely unexplored. \\

This gap is especially critical because ViTs typically require large, diverse datasets to perform effectively. The lack of sufficient training samples can lead to poor generalization, unstable convergence, and overfitting—issues that GANs are well-positioned to mitigate. However, little is known about how GAN-generated synthetic PSG images influence the attention mechanisms and feature extraction capabilities of ViTs. \\

Further research is needed to assess the quality, diversity, and impact of GAN-augmented samples on Transformer-based models. Exploring how GANs can be tailored to generate stage-specific, multi-channel sleep data in image format could unlock new opportunities for improving performance in low-data and class-imbalanced scenarios. \\
%--------------------------------------------------------------------------------------------------------

By addressing these gaps, this study aims to introduce a novel framework for fine-tuning pre-trained ViT models on PSG-based sleep stage classification, systematically evaluate different image transformation techniques, and explore GAN-based augmentation strategies tailored for Transformers.

\section{Summary of Methodological Alignment with Research Gaps}

The research methodology was carefully designed to directly address the gaps identified in the literature review (Section \ref{lit:5}). First, to resolve the limited application of pre-trained Vision Transformers (ViTs) in sleep stage classification, multiple experiments (Experiments 3, 4, and 6) were conducted to evaluate and fine-tune various ViT models, including ViT-B/16, ViT-B/32, DeiT-B, Swin-B, BEiT-B, and PVT-Small, on PSG-derived image representations. This approach specifically leverages transfer learning to enhance generalization and reduce computational overhead, as proposed in Gap 1 (Section \ref{lit:51}). Second, Experiment 1 systematically compares different image transformation techniques (Spectrograms, GADF, GASF, MTF) to identify the most informative representations for Transformer-based models, addressing Gap 2 (Section \ref{lit:52}) regarding the lack of consensus on optimal signal-to-image conversions. Third, to explore the underutilized potential of GAN-based augmentation in Transformer architectures (Gap 3, Section \ref{lit:53}), Experiments 5 and 6 employed a tuned DCGAN to generate synthetic PSG images, evaluating their impact on classification accuracy and data diversity across balanced, imbalanced, and mixed datasets. By aligning each methodological step with the findings and limitations of existing work, this study proposes a comprehensive and novel framework for advancing automated sleep stage classification using state-of-the-art deep learning techniques.